% Section 8: Conclusion (Target: 500 words)
\section{Conclusion}
\label{sec:conclusion}

Computational protein-protein interaction research is undergoing a fundamental transformation. This review has traced the trajectory from prediction---inferring whether, where, and how proteins interact---to design---creating proteins with desired interaction properties. This paradigm shift represents a qualitative change in how AI interfaces with biological discovery.

The technological foundations are established. Protein language models encode evolutionary and structural information at scale~\cite{rives2021esm1b,lin2023esm2}. Diffusion models generate novel structures conditioned on functional constraints~\cite{watson2023rfdiffusion,ingraham2023chroma}. Equivariant neural networks preserve geometric symmetries essential for accurate modeling~\cite{jing2021gvp,fuchs2020se3transformer}. Together, these enable design capabilities implausible just five years ago.

Yet significant challenges remain. The evaluation crisis suggests reported accuracies may be inflated~\cite{bernett2024bias}. Success rates vary widely (10--50\% across targets), with negative results systematically underreported. Experimental validation limits iterative improvement. Dynamic interactions, conditional specificity, and multi-protein complexes remain largely beyond current capabilities.

The implications extend beyond technical advances. Drug discovery transitions from high-throughput screening to targeted generation~\cite{zambaldi2024alphaproteo}. Synthetic biology can design interaction networks rather than repurposing natural components. Basic research gains hypothesis-driven experimental tools.

Perhaps most significantly, this transition reflects a philosophical shift: from AI as a tool for understanding what nature produced to AI as a partner in creating what nature has not. This does not diminish prediction---prediction remains the foundation upon which design builds~\cite{jumper2021alphafold,dauparas2022proteinmpnn}. Rather, it positions prediction as capability to be leveraged rather than endpoint in itself.

The coming years will determine whether current successes represent a sustained revolution or incremental advance peak. Progress depends on: better evaluation, faster validation, deeper understanding of designability, dynamics modeling, and interpretable methods. If addressed, the prediction-to-design transition will likely be remembered as a pivotal moment when computational biology substantially expanded its creative capabilities.

\vspace{1em}
\noindent
\textbf{Key Points:}
\begin{enumerate}
    \item AI-driven PPI research is transitioning from prediction to design---a qualitative expansion of capabilities enabled by protein language models, generative diffusion models, and equivariant neural networks.
    \item This constitutes a paradigm shift in the Kuhnian sense: reconceptualizing what questions are meaningful and what methods are appropriate.
    \item Design capabilities transform drug discovery, synthetic biology, and basic research, but success rates vary dramatically by target (10--50\%).
    \item Critical challenges---evaluation bias, validation bottlenecks, interpretability gaps, ethical considerations---must be addressed to realize the paradigm's potential.
    \item The future lies in prediction-design integration: AI systems iteratively creating, validating, and refining in closed-loop workflows.
\end{enumerate}

\vspace{1em}
\noindent
\textbf{Implications:}

\noindent
\textit{For Researchers:} Develop rigorous evaluation protocols, publish negative results, build computational-experimental feedback loops. Highest-impact work addresses where current methods fail: dynamic interactions, conditional specificity, multi-component assemblies.

\noindent
\textit{For Industry:} Prepare for accelerated discovery timelines and expanded druggable space. Investment in experimental validation infrastructure will differentiate successful programs. Economic advantage shifts from largest compound libraries to most efficient design-test cycles.

\noindent
\textit{For Funders:} Support fundamental challenges---better energy functions, dynamics modeling, interpretable design---rather than incremental benchmark improvements. Infrastructure for validation and data sharing amplifies methodological advances.

\vspace{1em}
\noindent
The prediction-to-design paradigm shift marks a transition from AI as passive observer of biology to AI as active creator of biological function. The technical foundations are in place; the challenges are known; the opportunity is unprecedented. Whether this transition fulfills its transformative promise will be determined by the collective efforts of researchers, developers, and practitioners in the coming years.
